% Parte do material do minicurso "Escrita de Artigos Científicos com LaTeX" (Conecta Hub)
% Exemplo 06 – Modelo de artigo científico
% Estrutura típica: resumo, introdução, método, resultados, conclusão e referências.
% Compile: pdflatex main.tex -> bibtex main -> pdflatex main.tex -> pdflatex main.tex

\documentclass[12pt, a4paper]{article}
\usepackage[utf8]{inputenc}
\usepackage[T1]{fontenc}
\usepackage[brazil]{babel}
\usepackage{amsmath, amssymb}
\usepackage{booktabs}
\usepackage{graphicx}
\usepackage[backend=bibtex, style=numeric]{biblatex}
\addbibresource{referencias.bib}   % copia de referencias.bib (pasta 05)
\usepackage{hyperref}                    % links clicáveis (use por último)

\title{Título do Artigo: Um Exemplo de Estrutura}
\author{Autor Um \and Autora Dois}
\date{\today}

\begin{document}

\maketitle

\begin{abstract}
Este resumo apresenta, em poucas linhas, o problema, o método, os principais
resultados e a conclusão do trabalho. Normalmente tem entre 150 e 250 palavras.
\end{abstract}

\noindent\textbf{Palavras-chave:} LaTeX; escrita científica; modelo.

\section{Introdução}
A introdução contextualiza o tema, apresenta o problema e enuncia a contribuição
do trabalho \cite{lamport1994}. Termine a introdução indicando a organização do texto.

\section{Trabalhos Relacionados}
Aqui você discute trabalhos anteriores \cite{knuth1984, exemplo2020}, destacando
o que já foi feito e qual é a lacuna que o seu trabalho preenche.

\section{Método}
Descreva o procedimento de forma reproduzível. Equações também fazem parte do método:
\begin{equation}
    y = \alpha x + \beta + \varepsilon
    \label{eq:modelo}
\end{equation}
onde $\alpha$ e $\beta$ são parâmetros estimados e $\varepsilon$ é o erro.

\section{Resultados}
A Tabela~\ref{tab:resultados} resume os resultados obtidos.

\begin{table}[h]
    \centering
    \caption{Resultados experimentais (valores ilustrativos).}
    \label{tab:resultados}
    \begin{tabular}{lcc}
        \toprule
        Método   & Precisão & Revocação \\
        \midrule
        Proposto & 0,92     & 0,88      \\
        Base     & 0,85     & 0,81      \\
        \bottomrule
    \end{tabular}
\end{table}

\section{Conclusão}
O método proposto superou a linha de base nas métricas avaliadas. Trabalhos futuros
incluem ampliar o conjunto de dados e comparar com outras abordagens.

\section*{Agradecimentos}
Opcional.

\printbibliography[title={Referências}]

\end{document}
